\subsection{应用：I. 连续实函数的向量空间}

设 X 是一个非空紧空间，H 是 Banach 空间 \(\mathcal{C}(\mathrm{X};\mathbf{R})\) 的一个线性子空间，它包含常数且分离 X 的点（GT, X, §4, No.~1, 定义 1）。我们给 \(\mathcal{H}\) 赋以由 \(\mathcal{C}(\mathrm{X};\mathbf{R})\) 的拓扑诱导的赋范空间拓扑，并用 \(\mathcal{H}'\) 表示这一赋范空间的对偶。对每个 \(x\in \mathrm{X}\)，映射 \(f\mapsto f(x)\) 是 \(\mathcal{H}\) 上的一个连续线性形式（\(\mathcal{H}\) 上 Dirac 测度 \(\varepsilon_{x}\) 的限制），因而是 \(\mathcal{H}'\) 中一个将记作 \(i_{\mathcal{H}}(x)\) 的元素，使得

\[
\langle f, i _ {\mathcal {H}} (x) \rangle = f (x)\tag{9}
\]

对每个函数 \(f \in \mathcal{H}\) 及每个 \(x \in X\) 成立。

映射 \(i_{\mathcal{H}}\)（从 X 到 \(\mathcal{H}'\) 中）当 \(\mathcal{H}'\) 赋以弱拓扑 \(\sigma(\mathcal{H}', \mathcal{H})\) 时是单射且连续的；第二个断言由定义及 (9) 立即得出；至于第一个，注意若 \(x, x'\) 是 X 的两个不同点，则由假设存在一个函数 \(h \in \mathcal{H}\) 使得 \(h(x) \neq h(x')\)，于是由 (9)，\(\langle h, i_{\mathcal{H}}(x) \rangle \neq \langle h, i_{\mathcal{H}}(x') \rangle\)，从而更有 \(i_{\mathcal{H}}(x) \neq i_{\mathcal{H}}(x')\)。因此象 \(i_{\mathcal{H}}(\mathrm{X})\) 是 \(\mathcal{H}'\) 的一个紧子集（对弱拓扑），且 \(i_{\mathcal{H}}\) 是 X 到 \(i_{\mathcal{H}}(\mathrm{X})\) 上的一个同胚。

命题 4. --- (i) \(i_{\mathcal{H}}(X)\) 在 \(\mathcal{H}'\) 中的闭凸包络 C（对弱拓扑 \(\sigma(\mathcal{H}', \mathcal{H})\)）是紧的。

\begin{enumerate}
\def\labelenumi{(\roman{enumi})}
\setcounter{enumi}{1}
\tightlist
\item
  为使点 \(i_{\mathcal{H}}(x)\) 是 \(C\) 的极值点，必须且只须：唯一的正测度 \(\lambda\)（\(X\) 上的）满足
\end{enumerate}

\[
h (x) = \int h d \lambda\tag{10}
\]

（对每个函数 \(h \in \mathcal{H}\)；这特别蕴涵 \(\lambda\) 的总质量为 1，因为 \(1 \in \mathcal{H}\)）就是 Dirac 测度 \(\varepsilon_x\)。

由此得出，对每个 \(h \in \mathcal{H}\)，函数 \(z' \mapsto \langle h, z' \rangle\)（\(C\) 上的）在 \(C\) 的至少一个极值点处达到其上确界（TVS, II, §7, No.~1, 命题 1），且这个点属于 \(i_{\mathcal{H}}(X)\)（同前，命题 2 的推论）。

\begin{enumerate}
\def\labelenumi{(\roman{enumi})}
\item
  由 (9)，\(\| i_{\mathcal{H}}(x)\| \leqslant 1\) 在赋范空间 \(\mathcal{H}'\) 中成立，换言之 \(i_{\mathcal{H}}(\mathrm{X})\) 是有界的，而这一断言由 \(\mathcal{H}'\) 赋以弱拓扑 \(\sigma (\mathcal{H}',\mathcal{H})\) 时是拟完备的这一事实得出（TVS, III, §4, No.~2, 定理 1 的推论 5）。
\item
  每个质量为 1 的正测度 \(\mu\)（在 \(i_{\mathcal{H}}(X)\) 上）都可借助同胚 \(i_{H}\) 通过结构转移而来自 X 上一个质量为 1 的正测度 \(\lambda\)，其中 Dirac 测度 \(\varepsilon_{i_{\mathcal{H}(x)}}\) 来自 \(\varepsilon_{x}\)。说 \(\mu\) 以 \(i_{\mathcal{H}(x)}\) 为重心，按定义即表示
\end{enumerate}

\[
\int_ {\mathrm{X}} \langle h, i _ {\mathcal {H}} (z) \rangle d \lambda (z) = \langle h, i _ {\mathcal {H}} (x) \rangle
\]

对每个函数 \(h \in H\) 成立。考虑到 (9)，断言 (ii) 无非是 No.~2 命题 3 的推论中关于 \(i_{\mathcal{H}}(x)\) 是 C 的极值点的判别准则的翻译。

我们把满足命题 4 的条件 (ii) 的点 \(x \in \mathbf{X}\) 称为 \(\mathcal{H}\)-极值点；用 \(\mathrm{Ch}_{\mathcal{H}}(\mathbf{X})\)（或简作 \(\mathrm{Ch}(\mathbf{X})\)）表示这些点所成之集，用 \(\check{\mathrm{S}}_{\mathcal{H}}(\mathbf{X})\)（或简作 \(\check{\mathrm{S}}(\mathbf{X})\)）表示 \(\mathrm{Ch}_{\mathcal{H}}(\mathbf{X})\) 在 \(\mathbf{X}\) 中的闭包。

命题 5. --- 每个函数 \(h \in \mathcal{H}\) 都在至少一个 \(\mathcal{H}\)-极值点处达到其上确界。

设 \(x\) 是 \(X\) 的一个点，\(h\) 是 \(\mathcal{H}\) 中的一个函数。关系 \(h(z) \leqslant h(x)\)（对一切 \(z \in X\)）可以写成 \(\langle h, i_{\mathcal{H}}(z) \rangle \leqslant \langle h, i_{\mathcal{H}}(x) \rangle\)（对一切 \(z \in X\)），因而它表示 \(\mathcal{H}'\) 中方程为 \(\langle h, t' \rangle = \langle h, i_{\mathcal{H}}(x) \rangle\) 的弱闭超平面是 \(i_{\mathcal{H}}(X)\) 的一个支撑超平面。已知（TVS, II, §7, No.~1, 命题 1 的推论）这样的超平面至少包含 \(i_{\mathcal{H}}(X)\) 的闭凸包络的一个极值点，而这样的点 \(i_{\mathcal{H}}(y)\) 按定义是一个 \(\mathcal{H}\)-极值点 \(y\) 的象；因此 \(h(y)\) 等于 \(h\) 在 \(X\) 中的上确界。

命题 6. --- 对每个点 \(x \in X\)，下列性质等价：

\begin{enumerate}
\def\labelenumi{\alph{enumi})}
\item
  \(x\) 是 \(\mathcal{H}\)-极值点。
\item
  对 x 在 X 中的每个开邻域 U 及每个 \(\varepsilon > 0\)，存在 H 中的一个函数 \(h \geqslant 0\)，使得 \(h(x) \leqslant \varepsilon\) 且 \(h(y) \geqslant 1\)（对每个 \(y \in X - U\)）。
\end{enumerate}

设 \(x\) 是 \(X\) 的任一点，\(f\) 是 \(\mathcal{C}(X; \mathbf{R})\) 中的一个函数；已知（TVS, II, §3, No.~1, 命题 1）诸数 \(\lambda(f)\)（对 \(X\) 上满足 \(\lambda(h) = h(x)\)——对每个函数 \(h \in \mathcal{H}\)——的一切正测度而取的）的下确界，等于诸数 \(h(x)\) 的上确界，其中 \(h\) 取遍函数 \(h \in \mathcal{H}\)（满足 \(h \leqslant f\)）所成之集。假定 \(x\) 是 \(\mathcal{H}\)-极值点；于是由命题 4, (ii)，对每个函数 \(f \in \mathcal{C}(X; \mathbf{R})\)，

\[
f (x) = \sup _ {h \in \mathscr {H}, h \leqslant f} h (x).\tag{11}
\]

为证明 \(a)\) 蕴涵 \(b)\)，我们取 \(f\) 为 X 到 \([0,1]\) 中的一个连续映射，其支集含于 U，且 \(f(x)=1\)；于是由 (11)，存在一个函数 \(h'\in\mathscr{H}\) 使得 \(h'\leqslant f\) 且 \(h'(x)\geqslant 1-\varepsilon\)。由于 \(1\in\mathscr{H}\)，函数 \(h=1-h'\) 满足 \(b)\) 的条件。

反之，假定条件 \(b)\) 成立；这一条件蕴涵 \(1 - h \leqslant \varphi_{\mathrm{U}}\)；于是对每个正测度 \(\lambda\)（\(X\) 上的，满足条件 (10)），我们有

\[
\lambda (\mathrm{U}) = \lambda (\varphi_ {\mathrm{U}}) \geqslant \lambda (1 - h) = 1 - h (x) \geqslant 1 - \varepsilon .
\]

由于按假设这一关系对每个 \(\varepsilon > 0\) 及 x 的每个开邻域 U 都成立，于是得到

\[
\lambda (\{x \}) = \inf _ {U} \lambda (U) \geqslant 1 - \varepsilon
\]

对每个 \(\varepsilon > 0\) 成立，因此 \(\lambda(\{x\}) = 1\)。由于 \(\lambda\) 是正的且总质量为 1，必有 \(\lambda = \varepsilon_x\)，这就证明了 \(x\) 是 \(\mathcal{H}\)-极值点，依据是命题 4, (ii)。

命题 7. --- 设 F 是 X 的一个闭子集。下列性质等价：

\begin{enumerate}
\def\labelenumi{\alph{enumi})}
\item
  F 包含 \(\check{\mathrm{S}}_{\mathcal{H}}(\mathrm{X})\)
\item
  对每个函数 \(h \in \mathcal{H}\)，集 \(F\) 都与 \(X\) 中使 \(h\) 达到其上确界的点所成之集相交。
\item
  对每个点 \(x \in X\)，存在一个总质量为 1 的正测度 \(\mu\)（在 \(X\) 上），使得 \(\operatorname{Supp}(\mu) \subset F\) 且 \(h(x) = \int h d\mu\)（对每个函数 \(h \in \mathcal{H}\)）。
\end{enumerate}

设 \(G = i_{\mathcal{H}}(F)\)。条件 \(a\)) 表示 \(G\) 包含 \(C\) 的极值点所成之集。条件 \(b\)) 表示 \(G\) 与 \(i_{\mathcal{H}}(X)\) 同 \(i_{\mathcal{H}}(X)\) 的每个闭支撑超平面的交相遇。最后，条件 \(c\)) 表示 \(i_{\mathcal{H}}(X)\) 的每个点都是某个支集含于 \(G\) 的测度的重心；由 No.~1 命题 1，这也等价于说 \(i_{\mathcal{H}}(X)\) 的闭凸包络等于 \(G\) 的闭凸包络。于是条件 \(a\))、\(b\)) 与 \(c\)) 的等价性由 TVS, II, §7, No.~1, 命题 2 的推论得出。

命题 8. --- 假定 X 是可度量化的。那么，X 的 H-极值点所成之集 \(\operatorname{Ch}_{\mathcal{H}}(\mathbf{X})\) 是 X 中一族可数开集的交，且对每个 \(x \in X\)，存在 X 上一个总质量为 1 的正测度 \(\mu\)，使得

\[
\mu (\mathrm{X} - \mathrm{Ch} _ {\mathcal {H}} (\mathrm{X})) = 0 \text { and } \int h d \mu = h (x)
\]

对每个 \(h \in \mathcal{H}\) 成立。

这就是 No.~2 的定理 1 借助同胚 \(x \mapsto i_{\mathcal{H}}(x)\) 通过结构转移得到的翻译，如同命题 5 中那样。

当把 H 换成一个由定义在 X 上、取值于 \(R \cup \{+\infty\}\) 的下半连续函数组成的集 P 时，本节的若干结果可以推广，这里假定 P 包含常数且满足 \(P + P \subset P\)（习题 2）。

*例. --- 取 X 为单位球 \(\| \mathbf{x} \| \leqslant 1\)（\(\mathbf{R}^3\) 中的），并设 \(\mathcal{H}\) 是 X 上连续函数的一个向量空间，它包含 \(\mathbf{R}^3\) 上仿射线性函数在 X 上的限制，且满足“极大值原理”，即对每个非常数函数 \(h \in \mathcal{H}\)，X 中使 \(h\) 达到其上确界的点所成之集含于球面 \(\mathbf{S}_2\)。于是由命题 5 与命题 7 容易得出 \(\mathrm{Ch}_{\mathcal{H}}(\mathrm{X}) = \check{\mathrm{S}}_{\mathcal{H}}(\mathrm{X}) = \mathbf{S}_2\)。满足前述条件的向量空间 \(\mathcal{H}\) 的一个重要例子是在 X 上连续且在开球 \(\| \mathbf{x} \| < 1\) 中调和的函数所成之集。对这些函数，人们证明：质量为 1 的正测度 \(\mu\)（满足 \(\operatorname{Supp}(\mu) \subset \mathbf{S}_2\) 且 \(h(\mathbf{x}) = \int h d\mu\)——对一切 \(h \in \mathcal{H}\)）当 \(\| \mathbf{x} \| < 1\) 时由 Poisson 公式给出

\[
d \mu (\mathbf {z}) = \frac {1 - \| \mathbf {z} \| ^ {2}}{\| \mathbf {z} - \mathbf {x} \| ^ {3}} d \sigma (\mathbf {z}),
\]

其中 \(\sigma\) 是 \(\mathbf{S}_2\) 上在正交群下不变且满足 \(\sigma(\mathbf{S}_2) = 1\) 的测度（Ch. VII, §3, 习题 8）。*

